% =====================================================================
%  Supplemental Material
% =====================================================================
\documentclass[aps,prl,onecolumn,superscriptaddress,amsmath,amssymb,floatfix]{revtex4-2}
\usepackage{graphicx,bm,xcolor,booktabs}
\usepackage[colorlinks=true,citecolor=blue!50!black,linkcolor=blue!50!black]{hyperref}
\newcommand{\gd}{\dot\gamma}
\newcommand{\ts}{\tau^{\star}}
\newcommand{\kT}{k_BT}
\newcommand{\todo}[1]{\textcolor{red}{[#1]}}
\renewcommand{\theequation}{S\arabic{equation}}
\renewcommand{\thetable}{S\arabic{table}}
\renewcommand{\thefigure}{S\arabic{figure}}

\begin{document}
\title{Supplemental Material: Unified kinetic description of flow stress, strain-rate sensitivity, dynamic strain aging, and creep in body-centered-cubic metals}
\author{Matthew Maron}
\author{Giacomo Po}
\author{Yang Li}
\maketitle

\section{S1. Components of the flow rule}

\subsection{Glide}
The mobile density moves with the kink-pair velocity
\begin{equation}
v_g=\frac{\ts b}{B_d}\exp\left\{-\frac{\Delta H_0}{2\kT}\max\left[\left(1-\left(\frac{\ts}{\tau_P}\right)^p\right)^q-\frac{T}{T_0},\,0\right]\right\}.
\end{equation}
For prismatic slip in Zr the drag coefficient interpolates smoothly between the kink drag $B_k$ and the free-flight drag $B_0$ as the barrier vanishes.

\subsection{Dislocation density}
The total density obeys (see Sec.~S5 for the origin of each term)
\begin{equation}
\frac{d\rho}{d\gamma}=k_1(\sqrt\rho+\Lambda_b)-k_2\rho-\frac{K D(T)}{\gd}\rho^2,
\end{equation}
with $k_2=k_{20}\exp(q_2/\kT)$ and $D=D_0\exp(-Q_D/\kT)$. The forest and mobile densities are $\rho_f=\beta\rho$ and $\rho_m=f_m\rho_f$. The flow stresses reported in the letter use the steady state, which is the root of
\begin{equation}
k_1(x+\Lambda_b)-k_2x^2-\frac{K D}{\gd}x^4=0,\qquad x=\sqrt\rho .
\end{equation}
The left-hand side is positive at $x=0$ and concave for $x>0$, so the positive root is unique. When the last term dominates, $x\propto(\gd/KD)^{1/3}$ and $\tau_a\propto(\gd/D)^{1/3}$. The boundary term is $\Lambda_b=\sum_ic_i/d_i$ over laths, blocks, prior-austenite grains and particles.

\subsection{Solute aging}
Arrested segments wait at forest junctions for $t_a=\rho_mb\lambda_f/\gd$ with $\lambda_f=\rho_f^{-1/2}$. The atmosphere concentration is
\begin{equation}
\frac{C}{C_0}=1+\sum_iC_i\left[1-e^{-x_i}\right],\qquad x_i=\left(t_a\nu_ie^{-Q_i/\kT}\right)^{\alpha_i},
\end{equation}
and the Taylor stress is multiplied by $s_c=\sqrt{C/C_0}$, so that $\tau_a=\alpha\mu b\,s_c(\sqrt{\rho_f}+\Lambda_b)$.

\subsection{Particle pinning}
Particles of areal density $\lambda_p^{-2}$ are met $n_p=\lambda_p^{-1}$ times per unit glide distance. The particle stress $\tau_p$ solves
\begin{equation}
\gd=\frac{\rho_mb}{n_p}\left(\nu_{\rm th}+\nu_{\rm cl}\right).
\end{equation}
Thermal detachment and glide-assisted local climb proceed at
\begin{align}
\nu_{\rm th}&=\nu_0\exp\left\{-\frac{F}{\kT}\left[1-\left(\frac{\tau_p}{\hat\tau}\right)^p\right]^q\right\},\\
\nu_{\rm cl}&=\frac{v_c}{h_{\rm eff}},\qquad
v_c=\frac{2\pi D_c}{b\ln(R/r_c)}\left[\exp\left(\frac{\tau_p\Omega}{\kT}\right)-1\right],\qquad
h_{\rm eff}=h\max\left(1-\frac{\tau_p}{\hat\tau},10^{-3}\right),
\end{align}
with $\hat\tau=\hat\tau_0\,\mu(T)/\mu_0$. The reduction of the climb height describes a line pressed against the particle, as in local-climb models~\cite{S_BrownHam}, and produces threshold-type creep. Both rates are non-decreasing in $\tau_p$, so $\tau_p$ is unique.

\subsection{Diffusional flow}
Nabarro--Herring and Coble flow add the rate
\begin{equation}
\gd_{\rm diff}=42\,\frac{\Omega\tau}{\kT d^2}\left(D_v+\frac{\pi\delta}{d}D_b\right).
\end{equation}

\subsection{Uniqueness and constant-stress evaluation}
With the state fixed by $(\gd,T)$, $v_g$ is strictly increasing in $\ts$ and $\tau_p$ is unique, so Eq.~(1) of the letter has exactly one solution for every imposed rate. The function $\tau(\gd)$ may decrease where aging gives $m<0$, but it is always single valued. For creep and for the maps, $\tau(\gd)$ is tabulated on $10^{-14}\le\gd\le10^{6}$~s$^{-1}$ at each temperature and inverted; where $\tau(\gd)$ is non-monotonic the lowest-rate branch that reaches the applied stress is taken, and $\gd_{\rm diff}$ is added.

\section{S2. Derivation of the rate-sensitivity decomposition}
Write $\tau=\ts+R$, where $\gd=G(\ts;\bm s)=\rho_mb\,v_g(\ts)$ is the intrinsic mobility and $R=\tau_a(\bm s)+\sum_g\tau_g$ collects all resistances. Along the solution,
\begin{equation}
d\ln\gd=A\,d\ts+\frac{\partial\ln G}{\partial\bm s}\cdot\frac{d\bm s}{d\ln\gd}\,d\ln\gd,\qquad A\equiv\frac{\partial\ln G}{\partial\ts},
\end{equation}
so that
\begin{equation}
\frac{d\tau}{d\ln\gd}=\frac1A+B,\qquad
B\equiv\frac{dR}{d\ln\gd}-\frac1A\,\frac{\partial\ln G}{\partial\bm s}\cdot\frac{d\bm s}{d\ln\gd}.
\end{equation}
With $\Lambda=AB$ this gives
\begin{equation}
m=\frac1\tau\frac{d\tau}{d\ln\gd}=\frac{1+\Lambda}{\tau A},\qquad V^\star=\frac{\kT}{m\tau}=\frac{\kT A}{1+\Lambda}.
\end{equation}
The mobility sensitivity is
\begin{equation}
A=\frac1\ts+\frac{\Delta H_0pq}{2\kT\tau_P}\,x^{p-1}(1-x^p)^{q-1},\qquad x=\frac{\ts}{\tau_P}.
\end{equation}
The contributions to $B$ are as follows. For the particle stress, $d\tau_p/d\ln\gd=1/A_p$ with $A_p=\sum_k(\nu_k/\nu)A_k$ and
\begin{align}
A_{\rm th}&=\frac{Fpq}{\kT\hat\tau}\,y^{p-1}(1-y^p)^{q-1},\qquad y=\frac{\tau_p}{\hat\tau},\\
A_{\rm cl}&=\frac{\Omega}{\kT}\frac{e^u}{e^u-1}+\frac{1}{\hat\tau-\tau_p},\qquad u=\frac{\tau_p\Omega}{\kT}.
\end{align}
Diffusion-controlled recovery lowers $\rho$ as $\gd$ decreases and therefore contributes $B>0$. Aging gives, with $t_a\propto\gd^{-1}$,
\begin{equation}
\frac{d\ln(C/C_0)}{d\ln\gd}=-\frac{\sum_iC_i\alpha_ix_ie^{-x_i}}{C/C_0}<0,
\end{equation}
so its contribution to $B$ is negative, and $m<0\iff\Lambda<-1$. In practice $A$ is evaluated analytically at the solved $\ts$, $d\tau/d\ln\gd$ is obtained by differentiating the full solution, and $B$ follows from its definition. The identity agrees with finite differences to better than $10^{-5}$ for all three materials, including points with $m<0$.

\section{S3. Material inputs and calibration}
Conventions: $\gd=\dot\epsilon/M_s$ and $\sigma=\tau/M_s$, with $M_s=0.333$ (W, Eurofer-97) and $0.45$ (Zr). The Zr and Eurofer-97 axial flow stresses are compared as $\sigma/3$. The global $m$ and $V^\star$ are least-squares fits over the experimental rate triplets.

The parameters were determined in the order of the hierarchy. The kink-pair parameters follow from the low-temperature, high-rate data; the athermal and storage parameters from the plateau; the aging parameters from the position of the stress hump and its shift with rate; and the recovery and climb parameters from the high-temperature, low-rate data. For W only level I parameters are required in the tensile range. For Zr the level II parameters were fitted to the resolved shear stresses of Derep \emph{et al.} at two rates, with a weak weight on $m$. For Eurofer-97 all parameters were fitted jointly to the flow stresses of Vanaja \emph{et al.} (three rates, 300--873~K) and to the minimum creep rates of Refs.~\cite{S_Fern,S_Yu}. The objective was the sum of squared relative stress errors plus $0.01$ times the sum of squared $\log_{10}$ rate errors, minimized by bounded Powell iterations. The activation energy for particle climb was set equal to that for forest recovery. The resulting mean flow-stress error is 4.1\%, and the creep rms error is 0.55 decades, with 52 of 54 rates within one decade. The diffusional-flow inputs were not fitted: $d=21~\mu$m (prior-austenite grain size), $D_v$ equal to the recovery diffusivity, and $D_b=2\times10^{-4}\exp(-1.8~\mathrm{eV}/\kT)$~m$^2$/s with $\delta=2b$. For W and Zr the recovery and diffusional terms used in the maps are based on self-diffusion data and are predictions.

\begin{table}[h]
\caption{Parameters. Values marked $^\dagger$ are lumped with geometry ($h$, $\ln R/r_c$) and are not bare physical constants.}
\begin{ruledtabular}
\begin{tabular}{llll}
 & W (I) & Zr (II) & Eurofer-97 (III)\\
\hline
$b$ (nm), $\mu_0$ (GPa), $T_m$ (K) & 0.2722, 161, 3695 & 0.3233, 33, 2128 & 0.2737, 80 [$\mu(T)$], 1811\\
$\Delta H_0$ (eV), $\tau_P$ (GPa), $p$, $q$ & 2.4, 0.8, 0.83, 1.69 & 3.75, 0.30, 0.86, 1.69 & 1.49, 0.68, 0.65, 1.7\\
$T_0/T_m$ & 0.9 & 0.7 & 0.9\\
$\alpha$ & 0.28 & 0.36 & 0.296\\
$k_1$ (m$^{-1}$), $k_{20}$, $q_2$ (eV) & $5.5\times10^9$, 500, $-0.011$ & $5.5\times10^9$, 500, $-0.005$ & $7.8\times10^{10}$, 6750, $-0.007$\\
$\beta$, $f_m$ & 1, 1 & 1, 1 & 0.1, 1\\
$\Lambda_b$ (m$^{-1}$) & -- & -- & $2.74\times10^6$\\
aging $C$, $Q$ (eV), $\alpha$ & -- & 0.30, 1.80, 2/3 (O) & 1.21, 0.97, 1/3 (C, N)\\
recovery $K$, $Q_D$ (eV) & $10^9$, 6.0 & $10^9$, 3.2 & $6.6\times10^{6}$, 2.92\\
particles $\hat\tau_0$ (MPa), $F$ (eV) & -- & -- & 188, 3.66\\
particle climb $D_{c,0}^\dagger$ (m$^2$/s), $Q$ (eV) & -- & -- & $5.4\times10^{-3}$, 2.92\\
\end{tabular}
\end{ruledtabular}
\end{table}

\section{S4. Data sources}
Tungsten: resolved shear stress, $m$ and $V^\star$ from micropillar compression, nanoindentation, indentation literature, IGW, W plate and arc-melted W \todo{references}. Zirconium: resolved shear stress and $V^\star$ from Derep \emph{et al.} \todo{ref}; $m$ from micro-cantilever tests on prismatic and basal slip \todo{Gong ref} and from Lee (1972, 2007) \todo{refs}. Eurofer-97 tensile data: Vanaja \emph{et al.}, J. Nucl. Mater. \textbf{424} (2012) \todo{pages}. Eurofer-97 creep: minimum creep rates digitized from Fig.~1 of Fern\'andez \emph{et al.}~\cite{S_Fern} (plate and bar, air and vacuum, 450--650~$^\circ$C, 53 points), together with the 550~$^\circ$C, 300~MPa steady-state rate of Yu \emph{et al.}~\cite{S_Yu}. The digitization uncertainty is about $\pm2\%$ in stress and $\pm0.05$ decade in rate. The digitized values are provided in \texttt{eurofer97\_creep\_data\_digitized.csv}.

\section{S5. Origin of the individual terms}
Table~\ref{tab:origin} lists where each ingredient of the flow rule comes from. The two recovery terms in the density evolution are physically distinct. The term $k_2\rho$ is dynamic recovery in the sense of Kocks and Mecking~\cite{S_Kocks1976,S_MK1981}: annihilation by cross slip and glide that happens per unit strain, so it does not depend explicitly on the strain rate (apart from the weak thermal activation of cross slip in $k_2$). The term $KD\rho^2/\gd$ is climb-controlled recovery. Its rate per unit time, $d\rho/dt=-KD\rho^2$, is the classical second-order law for the annihilation of dipoles and the coarsening of networks by climb~\cite{S_KMR1949,S_SL1975,S_S1977,S_Nes1995}. It is usually called static recovery because it proceeds without deformation, and it is converted to a strain basis by dividing by $\gd$. The same three-term form, $d\rho/d\epsilon=m/(bL)-\omega\rho-2\tau_LM_{\rm climb}\rho^2/\dot\epsilon$ with a climb mobility proportional to the self-diffusivity, is used by Sandstr\"om and co-workers to describe stress--strain curves and creep of copper with one set of equations~\cite{S_SL1975,S_SH2012}. Related forms appear in the internal-variable models of Nes~\cite{S_Nes1997} and of Roters, Raabe and Gottstein~\cite{S_RRG2000}. Recent mesoscale simulations~\cite{S_KC2022} indicate that the $\rho^2$ law is an approximation to a more complex climb-recovery kinetics.

\begin{table}[h]
\caption{Origin of each term.}
\label{tab:origin}
\begin{ruledtabular}
\begin{tabular}{lll}
Term & Form used & Origin\\
\hline
Orowan relation & $\gd=\rho_mbv_g$ & \cite{S_Orowan1940}\\
Kink-pair mobility & Eq.~(S1) & \cite{S_Po2016,S_KAA1975,S_DR1964}\\
Stress superposition & $\tau=\tau_a+\ts+\sum_g\tau_g$ & \cite{S_Seeger1958,S_KAA1975}\\
Taylor hardening & $\alpha\mu b\sqrt{\rho_f}$ & \cite{S_Taylor1934,S_KM2003}\\
Forest storage & $k_1\sqrt\rho$ & \cite{S_Kocks1976,S_MK1981}\\
Boundary storage & $k_1\Lambda_b$ & \cite{S_EM1984,S_Estrin1996}\\
Dynamic recovery & $k_2\rho$, $k_2=k_{20}e^{q_2/\kT}$ & \cite{S_Kocks1976,S_MK1981,S_KM2003}\\
Climb (static) recovery & $KD\rho^2/\gd$ & \cite{S_KMR1949,S_SL1975,S_S1977,S_Nes1995,S_SH2012}\\
Waiting time & $t_a=\rho_mb\lambda_f/\gd$ & \cite{S_McC1988,S_KE1990}\\
Atmosphere kinetics & Eq.~(S4) & \cite{S_CB1949,S_Louat1981,S_Cheng2001}\\
Solute strengthening & $s_c=\sqrt{C/C_0}$ & \cite{S_Cheng2001,S_MuK1979}\\
Thermal detachment & $\nu_{\rm th}$, Eq.~(S6) & \cite{S_KAA1975,S_RA1990}\\
Climb velocity & $v_c$, Eq.~(S7) & \cite{S_HL1982}\\
Local climb, threshold & $h_{\rm eff}=h(1-\tau_p/\hat\tau)$ & \cite{S_BrownHam,S_AA1982}\\
Diffusional flow & Eq.~(S8) & \cite{S_Nabarro1948,S_Herring1950,S_Coble1963,S_FA1982}\\
Natural creep law & $\tau_a\propto(\gd/D)^{1/3}$ & \cite{S_Poirier1985,S_Blum2002}\\
\end{tabular}
\end{ruledtabular}
\end{table}

\begin{thebibliography}{99}
\bibitem{S_Orowan1940} E. Orowan, Proc. Phys. Soc. \textbf{52}, 8 (1940).
\bibitem{S_Po2016} G. Po, Y. Cui, D. Rivera, D. Cereceda, T. D. Swinburne, J. Marian, and N. Ghoniem, Acta Mater. \textbf{119}, 123 (2016).
\bibitem{S_KAA1975} U. F. Kocks, A. S. Argon, and M. F. Ashby, Prog. Mater. Sci. \textbf{19}, 1 (1975).
\bibitem{S_DR1964} J. E. Dorn and S. Rajnak, Trans. Metall. Soc. AIME \textbf{230}, 1052 (1964).
\bibitem{S_Seeger1958} A. Seeger, in \emph{Handbuch der Physik}, Vol.~VII/2, edited by S. Fl\"ugge (Springer, Berlin, 1958).
\bibitem{S_Taylor1934} G. I. Taylor, Proc. R. Soc. London A \textbf{145}, 362 (1934).
\bibitem{S_KM2003} U. F. Kocks and H. Mecking, Prog. Mater. Sci. \textbf{48}, 171 (2003).
\bibitem{S_Kocks1976} U. F. Kocks, J. Eng. Mater. Technol. \textbf{98}, 76 (1976).
\bibitem{S_MK1981} H. Mecking and U. F. Kocks, Acta Metall. \textbf{29}, 1865 (1981).
\bibitem{S_EM1984} Y. Estrin and H. Mecking, Acta Metall. \textbf{32}, 57 (1984).
\bibitem{S_Estrin1996} Y. Estrin, in \emph{Unified Constitutive Laws of Plastic Deformation}, edited by A. S. Krausz and K. Krausz (Academic Press, San Diego, 1996), p.~69.
\bibitem{S_KMR1949} D. Kuhlmann, G. Masing, and J. Raffelsieper, Z. Metallkd. \textbf{40}, 241 (1949).
\bibitem{S_SL1975} R. Sandstr\"om and R. Lagneborg, Acta Metall. \textbf{23}, 387 (1975).
\bibitem{S_S1977} R. Sandstr\"om, Acta Metall. \textbf{25}, 897 (1977).
\bibitem{S_Nes1995} E. Nes, Acta Metall. Mater. \textbf{43}, 2189 (1995).
\bibitem{S_Nes1997} E. Nes, Prog. Mater. Sci. \textbf{41}, 129 (1997).
\bibitem{S_SH2012} R. Sandstr\"om and J. Hallgren, J. Nucl. Mater. \textbf{422}, 51 (2012).
\bibitem{S_RRG2000} F. Roters, D. Raabe, and G. Gottstein, Acta Mater. \textbf{48}, 4181 (2000).
\bibitem{S_KC2022} A. A. Kohnert and L. Capolungo, npj Comput. Mater. \textbf{8}, 104 (2022).
\bibitem{S_McC1988} P. G. McCormick, Acta Metall. \textbf{36}, 3061 (1988).
\bibitem{S_KE1990} L. P. Kubin and Y. Estrin, Acta Metall. Mater. \textbf{38}, 697 (1990).
\bibitem{S_CB1949} A. H. Cottrell and B. A. Bilby, Proc. Phys. Soc. A \textbf{62}, 49 (1949).
\bibitem{S_Louat1981} N. Louat, Scr. Metall. \textbf{15}, 1167 (1981).
\bibitem{S_Cheng2001} J. Cheng, S. Nemat-Nasser, and W. Guo, Mech. Mater. \textbf{33}, 603 (2001).
\bibitem{S_MuK1979} R. A. Mulford and U. F. Kocks, Acta Metall. \textbf{27}, 1125 (1979).
\bibitem{S_RA1990} J. R\"osler and E. Arzt, Acta Metall. Mater. \textbf{38}, 671 (1990).
\bibitem{S_HL1982} J. P. Hirth and J. Lothe, \emph{Theory of Dislocations}, 2nd ed. (Wiley, New York, 1982).
\bibitem{S_AA1982} E. Arzt and M. F. Ashby, Scr. Metall. \textbf{16}, 1285 (1982).
\bibitem{S_Nabarro1948} F. R. N. Nabarro, in \emph{Report of a Conference on the Strength of Solids} (Physical Society, London, 1948), p.~75.
\bibitem{S_Herring1950} C. Herring, J. Appl. Phys. \textbf{21}, 437 (1950).
\bibitem{S_Coble1963} R. L. Coble, J. Appl. Phys. \textbf{34}, 1679 (1963).
\bibitem{S_FA1982} H. J. Frost and M. F. Ashby, \emph{Deformation-Mechanism Maps} (Pergamon, Oxford, 1982).
\bibitem{S_Poirier1985} J.-P. Poirier, \emph{Creep of Crystals} (Cambridge University Press, Cambridge, 1985).
\bibitem{S_Blum2002} W. Blum, P. Eisenlohr, and F. Breutinger, Metall. Mater. Trans. A \textbf{33}, 291 (2002).
\bibitem{S_BrownHam} L. M. Brown and R. K. Ham, in \emph{Strengthening Methods in Crystals}, edited by A. Kelly and R. B. Nicholson (Elsevier, Amsterdam, 1971), p.~9.
\bibitem{S_Fern} P. Fern\'andez, A. M. Lancha, J. Lape\~na, R. Lindau, M. Rieth, and M. Schirra, Fusion Eng. Des. \textbf{75--79}, 1003 (2005).
\bibitem{S_Yu} G. Yu, N. Nita, and N. Baluc, Fusion Eng. Des. \textbf{75--79}, 1037 (2005).
\end{thebibliography}
\end{document}
